\documentclass[11pt,a4paper]{article}
\usepackage{fontspec}
\setmainfont{Times New Roman}
\setmonofont[Scale=.85]{Consolas}
\usepackage{amsmath,amssymb,booktabs,array,longtable,xcolor,geometry,fancyhdr,hyperref}
\geometry{margin=20mm,headheight=15pt}
\hypersetup{unicode=true,colorlinks=true,linkcolor=teal!60!black,urlcolor=blue!60!black,pdftitle={Tiền xử lý dữ liệu PPS bằng Excel}}
\setlength{\parindent}{0pt}
\setlength{\parskip}{5pt}
\setlength{\emergencystretch}{2em}
\renewcommand{\arraystretch}{1.2}
\allowdisplaybreaks
\pagestyle{fancy}\fancyhf{}\fancyhead[L]{TIỀN XỬ LÝ DỮ LIỆU PPS BẰNG EXCEL}\fancyfoot[R]{\thepage}
\newcommand{\xl}[1]{{\ttfamily\small\detokenize{#1}}}
\begin{document}
\begin{center}
{\LARGE\bfseries Tiền xử lý dữ liệu bằng Excel}\par\medskip
{\large Làm trước khi nhập vào tool; luôn có đường đổi về bài toán gốc}
\end{center}
Tài liệu riêng này dùng các phép tính Excel cơ bản. Đi kèm là
\texttt{Tien\_xu\_ly\_Excel.xlsx}, gồm dữ liệu tự dựng và các công thức có thể kéo xuống.
Giữ sheet dữ liệu gốc nguyên trạng. Các biến đổi không được xem là bước ngầm của bộ giải.

\tableofcontents
\section{Quy trình sáu bước}
\begin{enumerate}
\item Đọc đề: tìm giá trị hàm, đạo hàm, nghiệm ngược, mô hình thực nghiệm hay tích phân? Đề bắt buộc phương pháp nào, bao nhiêu mốc và điều kiện biên nào?
\item Giữ cột ID, $x_i,y_i$ gốc; ghi số quan sát $N$, đơn vị, hai giá trị đích và quy ước làm tròn. Với $N$ dòng, chỉ số là $0,\ldots,N-1$.
\item Kiểm tra ô số, ô thiếu, dấu, mốc trùng, khoảng cách và nhánh. Sắp \emph{cả bảng}, không chỉ riêng cột $x$.
\item Chọn phép xử lý có lý do; ghi công thức thuận, công thức ngược và đại lượng bị thay đổi. Lọc quan sát là một quyết định riêng, không phải phép tương đương đại số.
\item Nhập đúng dữ liệu và mục tiêu vào tool; kiểm tra bảng mốc/cơ sở/điều kiện biên của lượt chạy trước khi dùng kết quả.
\item Đổi về biến gốc; thế lại dữ liệu gốc. Khi khớp, ghi sai số được tính trên tập nào và trên thang nào.
\end{enumerate}

\section{Bố trí và kiểm tra bảng gốc}
Giả sử hàng 1 là tiêu đề, cột A là ID, B là $x$, C là $y$; dữ liệu ở hàng 2 đến 105. Cần sửa số hàng cuối theo dữ liệu thực.
\begin{center}\small
\begin{tabular}{p{.27\linewidth}p{.66\linewidth}}
\toprule Kiểm tra & Công thức tại hàng 2 hoặc công thức tổng hợp\\\midrule
Hai ô đều là số & \xl{=AND(ISNUMBER(B2),ISNUMBER(C2))}\\
Số dòng có $x$ dạng số & \xl{=COUNT(B2:B105)}; phải đối chiếu cả cột $y$\\
Mốc trùng & \xl{=COUNTIF($B$2:$B$105,B2)>1}\\
Bước kế tiếp & \xl{=B3-B2}, kéo đến khoảng cuối\\
Đều theo bước gốc & \xl{=ABS((B3-B2)-$J$2)<=1E-8*ABS($J$2)}\\
Số tung độ âm/dương/0 & \xl{=COUNTIF(C2:C105,"<0")}, đổi điều kiện thành \xl{">0"} hoặc \xl{"=0"}\\
Đổi dấu quanh mức $y^*$ & \xl{=(C2-$J$3)*(C3-$J$3)<0}\\
Trùng đúng mức & \xl{=C2=$J$3}; kiểm tra riêng cả mốc cuối\\
\bottomrule
\end{tabular}\end{center}
Đặt J2 là bước tham chiếu $h$, J3 là mức $y^*$. Công thức dùng tên hàm Excel tiếng Anh và dấu phẩy; máy dùng dấu chấm phẩy thì đổi dấu phân cách đối số, không đổi dấu thập phân của ô số.

Vào \textbf{Data $\to$ Sort}, chọn toàn vùng A:C và sắp cột B tăng dần.
Nếu Excel hỏi mở rộng vùng, chọn \emph{Expand the selection}. Hai cột $x,y$ phải luôn đi cùng nhau.
Không dùng \emph{Remove Duplicates} tùy tiện: nội suy/spline cần xử lý mốc trùng, nhưng bình phương tối thiểu có thể giữ các phép đo tại cùng $x$.

\section{Tịnh tiến và đổi đơn vị: phải đổi cả mục tiêu}
\subsection{Công thức chung}
Chọn $s\ne0,r\ne0$; để thao tác và thứ tự mốc dễ theo dõi thường chọn $s,r>0$:
\[
u=\frac{x-c}{s},\qquad v=\frac{y-d}{r},\qquad
x=c+su,\qquad y=d+rv.
\]
Tịnh tiến thuần túy là $s=r=1$. Đổi gốc $x$ giúp xử lý năm/thời gian hoặc tọa độ lớn; đổi đơn vị giúp các số có cỡ hợp lý.
Không thể lấy lại chữ số đã mất trong tệp gốc hay trong ô Excel vốn chỉ giữ khoảng 15 chữ số có nghĩa.

Đặt J2=$c$, J3=$s$, J4=$d$, J5=$r$ trong sheet biến đổi:
\begin{center}
D2: \xl{=(B2-$J$2)/$J$3},\quad
E2: \xl{=(C2-$J$4)/$J$5}.
\end{center}
Kéo xuống đến dòng cuối. Điểm cần tính và mức ngược phải đổi theo cùng công thức:
\[
\bar u=(\bar x-c)/s,\qquad \bar v=(\bar y-d)/r.
\]
Ví dụ $x=1000,1002,1004$; $y=5,9,17$. Chọn $c=1000,s=2,d=5,r=4$,
được $(u,v)=(0,0),(1,1),(2,3)$. Muốn tính ở $x=1003$ phải nhập $u=1.5$.
Nếu tool trả $V(1.5)=1.875$, đổi lại $y=5+4(1.875)=12.5$.
Muốn tìm $x$ ứng với $y=12.5$ phải nhập mức $v=1.875$; từ nghiệm $u$ đổi lại $x=1000+2u$.

\subsection{Nội suy, spline và đạo hàm}
Đa thức/mảnh spline tương ứng có quan hệ
\[
P(x)=d+rV\left(\frac{x-c}{s}\right),\qquad
P^{(k)}(x)=\frac{r}{s^k}V^{(k)}(u)\quad(k\ge1).
\]
Nếu đề cho đạo hàm ở biên, dữ liệu biên nhập vào bài theo $u,v$ là
\[
V^{(k)}(u_\text{biên})=\frac{s^k}{r}P^{(k)}(x_\text{biên}).
\]
Với ví dụ trên, nếu $V'(1.5)=2$ thì $P'(1003)=(4/2)2=4$.
Đổi riêng $x$ mà giữ nguyên trị số của đạo hàm biên sẽ tạo một bài spline khác.
Điều kiện tự nhiên (đạo hàm bậc hai bằng 0) vẫn là 0; điều kiện khác 0 phải đổi thang.
Horner không tự biết đa thức đang theo $u$ hay $x$: phải nhập hệ số theo đúng biến và điểm thế tương ứng.

\subsection{Tích phân}
Với $u_a=(a-c)/s,u_b=(b-c)/s$:
\[
\int_a^b P(x)\,dx=s\left[d(u_b-u_a)+r\int_{u_a}^{u_b}V(u)\,du\right].
\]
Không chỉ nhân $rs$ khi đã trừ một mức nền $d\ne0$.
Cận đạo hàm cũng đổi: $M_{k,V}=|s|^kM_{k,P}/|r|$ nếu dùng để chọn lưới theo cận lý thuyết.

\subsection{Phương trình vi phân và biên}
Cho $t=c+s\tau$, $y=d+rv$, phương trình $dy/dt=f(t,y)$ trở thành
\[
\frac{dv}{d\tau}=\frac{s}{r}f(c+s\tau,d+rv).
\]
Đổi cả khoảng thời gian, bước, điều kiện đầu và biểu thức vế phải. Với hệ nhiều thành phần, mỗi thành phần có thể có mức nền và thang riêng.

Điều kiện $\alpha y'(a)+\beta y(a)=\gamma$ đổi thành
\[
\frac{\alpha r}{s}v'(u_a)+\beta r v(u_a)=\gamma-\beta d.
\]
Với bài $-(p y')'+qy=f$, đặt $\widetilde p(u)=p(c+su)$ và tương tự cho $q,f$:
\[
-(\widetilde p v')'+s^2\widetilde q v
=\frac{s^2}{r}(\widetilde f-\widetilde q d).
\]
Đây là bài mới phải nhập đủ các hệ số đã đổi, không chỉ thay bảng mốc.
Với trị riêng và biên đồng nhất, giữ $d=0$; tịnh tiến tung độ khác 0 nói chung phá dạng bài trị riêng đồng nhất.

\section{Logarit và dấu: quyết định trước khi chạy khớp}
\subsection{Mô hình hàm mũ có biên độ âm}
Với $y=Ae^{bx+c/x}$ và $x_i\ne0$, dấu $y$ luôn là dấu $A$.
Nếu \emph{toàn bộ} $y_i<0$:
\[
-y=(-A)e^{bx+c/x}
\ \Longrightarrow\ \ln(-y)=\ln(-A)+bx+\frac cx.
\]
Đặt $Z=\ln(-y)$, $C_0=\ln(-A)$; cơ sở $1,x,1/x$.
Trong Excel dùng \xl{=IF(C2<0,LN(-C2),NA())} để hiện lỗi miền thay vì giấu dữ liệu.
Trong tool chọn \textbf{Hàm mũ}, nhập cơ sở \xl{x;1/x}, và vẫn đưa \emph{dữ liệu $x,y$ gốc} vào; tool tự lấy $\ln|y|$ và thêm cột 1.
Nếu đã đưa cột log vào một bài \textbf{Tuyến tính theo cơ sở} thì phải khôi phục $A=-e^{C_0}$ thủ công; không chọn Hàm mũ thêm lần nữa để tránh log hai lần.

Nếu mọi $y_i>0$, thay $\ln(-y)$ bằng $\ln y$, $A=+e^{C_0}$.
Nếu $y$ lẫn dấu hoặc có 0, phép tuyến tính hóa log không dùng được cho toàn bảng.
\textbf{Không âm thầm đổi dấu từng điểm bằng ABS rồi khớp như cùng một hàm.}

\subsection{Loại điểm và đặt trọng số 0}
Trong mẫu người dùng gửi, có 104 điểm và hai điểm cuối dương; bài mẫu chỉ khớp 102 điểm âm.
Đó là \textbf{một lựa chọn tập dữ liệu khác}, không suy ra được chỉ từ việc muốn lấy log.
Nếu đề cho phép loại: ghi ID, lý do, số dòng trước/sau, mô hình trước/sau và mẫu tính sai số.
Với bộ log của tool, đặt trọng số 0 vẫn không làm một quan sát sai miền trở thành hợp lệ: kiểm tra miền chạy trên các quan sát đã nhập.

Nếu cần giữ bảng lẫn dấu, nhánh \textbf{Hàm mũ — SSE thang y} có thể tối ưu trực tiếp trên $y$ với tham số đầu phù hợp.
Mô hình vẫn có một dấu; nó có thể xấp xỉ dữ liệu lẫn dấu bằng phần dư, nhưng không đi qua tất cả các điểm và không tối ưu cùng mục tiêu với bài log đã lọc.

\subsection{Không tịnh tiến để ``cứu'' log rồi gọi là cùng mô hình}
Nếu đặt $v=y+K>0$ rồi khớp $v=Ae^{bx}$, hàm khôi phục là
\[
y=Ae^{bx}-K,
\]
khác họ $y=Ae^{bx}$ ban đầu. Chỉ làm khi đề cho họ có mức nền tương ứng.
Với mô hình $A x^b$, tịnh tiến $x$ cho $(x-c)^b$, nói chung không còn là $x^b$.
Với $Ae^{bx+c/x}$, tịnh tiến $x=c_0+su$ buộc cơ sở là $c_0+su$ và $1/(c_0+su)$; không được tự thay bằng $u,1/u$.
Với hàm mũ đơn $Ae^{bx}$, đổi $x=c_0+su$ vẫn cùng họ:
\[
y=\underbrace{Ae^{bc_0}}_{\widetilde A}e^{(bs)u},\qquad
b=\widetilde b/s,\quad A=\widetilde A e^{-\widetilde b c_0/s}.
\]

\subsection{Bảng quyết định theo từng mô hình khớp}
\begin{center}\small
\begin{tabular}{p{.27\linewidth}p{.66\linewidth}}\toprule
Họ ban đầu & Điều cần giữ khi tiền xử lý\\\midrule
Tuyến tính theo cơ sở & Đổi $x=c+su$ thì đổi mỗi cơ sở thành $\varphi_j(c+su)$. Trừ mức nền của $y$ cần biểu diễn lại phần hằng; mô hình không có cơ sở 1 có thể đổi họ.\\
Đa thức đủ bậc 0 đến $p$ & Tịnh tiến/đổi thang $x,y$ vẫn cùng họ đủ bậc; phải khai triển lại tham số khi viết theo $x$.\\
$A x^b$ & Với $x=su$, $s>0$, biên độ theo $u$ là $\widetilde A=A s^b$, nên $A=\widetilde A/s^b$. Tịnh tiến $x$ hoặc trừ mức nền $y$ nói chung đổi họ.\\
$A B^x$ & Với $x=c+su$, $\widetilde A=A B^c$, $\widetilde B=B^s$; khôi phục $B=\widetilde B^{1/s}$, $A=\widetilde A/B^c$.\\
$A\exp(\sum b_j\varphi_j)$ & Đổi từng cơ sở theo biến mới; cột hằng trong số mũ biểu diễn log của biên độ. Không trừ mức nền $y$ nếu họ không chứa mức nền.\\
$\ln(\sum c_j\varphi_j)$ & Cần tổng bên trong log dương ở điểm đánh giá. Trừ $d$ trên $y$ tương ứng chia tổng bên trong cho $e^d$; nhân $y$ bởi một hệ số khác 1 nói chung đổi dạng tổ hợp cơ sở.\\
$1/(a+bx)$ & Với $x=c+su$, mẫu số là $(a+bc)+(bs)u$. Tịnh tiến $y$ nói chung không còn dạng nghịch đảo đường thẳng này.\\
$x/(a+bx)$ & Tịnh tiến $x=c+su$ đổi cả tử và mẫu; không thay tử bằng $u$ một cách ngầm định.\\
$a+b\ln x$ & Đổi $x=su$, $s>0$, thành $(a+b\ln s)+b\ln u$. Tịnh tiến $x$ khác 0 cho $\ln(c+su)$, không phải $\ln u$.\\
Hai nhánh SSE thang $y$ & Giữ hàm gốc và thang mục tiêu; tham số đầu đổi cùng biến. Khớp log chỉ là một cách khởi tạo khi đủ miền, không thay mục tiêu LM.\\\bottomrule
\end{tabular}\end{center}
Với cùng một phép đổi $v=(y-d)/r$ cho mọi dòng và họ mô hình đã đổi đúng,
SSE trên $v$ bằng SSE trên $y$ chia $r^2$, nên có cùng nghiệm cực tiểu sau ánh xạ về.
Điều này không áp dụng cho log, nghịch đảo hoặc những phép đổi phi tuyến của tung độ.

\section{Lập bảng tổng bằng Excel}
Ví dụ đã giữ hợp lệ $N$ dòng cho mô hình $Z=C_0+bx+c/x$.
Đặt B=$x$, C=$y$, D=$Z=\ln|y|$, E=$x^2$, F=$1/x$, G=$1/x^2$, H=$xZ$, I=$Z/x$.
\begin{center}\small
\begin{tabular}{llll}
\toprule D2 & E2 & F2 & G2\\\midrule
\xl{=LN(ABS(C2))}&\xl{=B2^2}&\xl{=1/B2}&\xl{=1/B2^2}\\\bottomrule
\end{tabular}\par\smallskip
H2: \xl{=B2*D2},\quad I2: \xl{=D2/B2}.
\end{center}
Dùng \xl{=SUM(B2:B103)}, \xl{=SUM(E2:E103)}, v.v. nếu có đúng 102 dòng dữ liệu.
Viết vào bài thi:
\[
\sum x_i=x_0+x_1+\cdots+x_{N-1},\quad
\sum x_iZ_i=x_0Z_0+x_1Z_1+\cdots+x_{N-1}Z_{N-1}.
\]
\[
\begin{bmatrix}N&\sum x_i&\sum1/x_i\\
\sum x_i&\sum x_i^2&N\\
\sum1/x_i&N&\sum1/x_i^2\end{bmatrix}
\begin{bmatrix}C_0\\b\\c\end{bmatrix}=
\begin{bmatrix}\sum Z_i\\\sum x_iZ_i\\\sum Z_i/x_i\end{bmatrix}.
\]
Hai ô chéo là $N$ vì $x_i(1/x_i)=1$ trên từng dòng. Có trọng số thì các ô này là $\sum w_i$, không phải số dòng.

Sau khi có tham số, thêm cột $\widehat y=A\exp(bx+c/x)$ và $e=y-\widehat y$:
\[
\operatorname{RMSE}_y=\sqrt{\frac{e_0^2+\cdots+e_{N-1}^2}{N}}.
\]
Excel có thể dùng \xl{=SQRT(SUMSQ(L2:L103)/COUNT(L2:L103))} với cột L là phần dư.
Đây là RMSE trên $y$; mục tiêu được tối thiểu hóa bằng log là tổng bình phương trên $Z$, hai đại lượng không giống nhau.
Không thay mẫu số bằng $N-p$ rồi vẫn gọi là RMSE; đó là một định nghĩa thống kê khác.

\section{Chọn mốc và nhánh trong Excel}
\textbf{Nội suy thuận.} Tìm mốc gần $\bar x$, rồi chọn đủ mốc theo đề.
Stirling cần số mốc lẻ và đối xứng quanh mốc giữa; Bessel dùng số mốc chẵn.
Không chọn một tập lỗ chỗ rồi coi vẫn có bước lưới cũ.
Ghi bảng ngang: chỉ số cục bộ, ID nguồn, $x_j,y_j$; số mốc quyết định bậc tối đa.

\textbf{Nội suy ngược.} Quét các mốc đúng mức và các cặp đổi dấu quanh $\bar y$.
Kiểm tra chiều tăng/giảm của $y$ trên từng nhánh; cùng mức có thể có nhiều nghiệm.
Không sắp cột $y$ toàn cục độc lập để biến dữ liệu nhiều nhánh thành một hàm ngược.
Trong tool, khoảng nghiệm không thay thế tập mốc: chọn ID hoặc nhập riêng các cặp muốn dùng.
Lặp tiến/lùi vẫn lấy gốc đầu/cuối tập đã chọn, không tự đổi sang gốc ngay cạnh nghiệm.

Nếu không đổi dấu, chỉ kết luận \emph{phép quét không tìm thấy khoảng đổi dấu}; nghiệm tiếp xúc giữa các mốc vẫn có thể tồn tại.
Nếu hai mốc liên tiếp bằng mức đích, dữ liệu chỉ cho thấy một đoạn phẳng trên mẫu, chưa chứng minh hàm thật hằng trên đoạn đó.

\section{Nhập vào tool và kiểm tra ngược}
\begin{enumerate}
\item Lưu vùng dùng để nhập thành hai cột đầu $x,y$. Workbook kèm có ba sheet nhập dữ liệu:\par
\texttt{Nhap\_tool\_goc}, \texttt{Nhap\_tool\_doi}, \texttt{Nhap\_tool\_log\_goc}.\par
Các sheet này tránh nhập nhầm ID thành $x$. Nếu tool đọc file Excel, ID và các cột phụ đặt sau hai cột này hoặc giữ ở sheet khác.
\item Nếu biến đổi, đặt tên sheet rõ ràng, ghi $c,s,d,r$, điểm/mức đích mới và đơn vị. Khi dùng một mô hình có phép biến đổi bên trong, nhập dữ liệu gốc đúng như hướng dẫn của mô hình.
\item Mở kết quả và đọc tập mốc thực dùng, cơ sở, gốc, điều kiện biên, bộ giải, thang mục tiêu và tiêu chí dừng.
\item Đổi kết quả về biến gốc; kiểm tra vài mốc gốc, đạo hàm biên gốc và phần dư ở mức đích gốc.
\item Ghi rõ mọi dòng loại, mọi đổi đơn vị và sai số trên đúng tập. Giữ đủ chữ số trung gian; làm tròn ở phần trình bày cuối.
\end{enumerate}
Các công thức trong workbook là công thức Excel, không có giá trị cache giả. Excel sẽ tính lại khi mở; đọc bằng thư viện với chế độ chỉ lấy giá trị trước khi Excel lưu lại có thể thấy ô công thức trống.
\end{document}
